
\documentclass[runningheads]{llncs}

\usepackage[T1]{fontenc}
\usepackage{graphicx}
\usepackage{array}
\usepackage{ragged2e}
\usepackage{physics}
\usepackage{amsmath}
\usepackage{amssymb}


\begin{document}
%
\title{Auto Quantum Machine Learning for Multisource Classification}

\author{Tomasz Rybotycki\inst{1, 2, 3}\orcidID{0000-0003-2493-0459} \and
	Sebastian Dziura\inst{1}\orcidID{0009-0008-6786-511X} \and
	Piotr Gawron\inst{1, 2}\orcidID{0000-0001-7476-9160}}
%
\authorrunning{T. Rybotycki, S. Dziura, P. Gawron}
% First names are abbreviated in the running head.
% If there are more than two authors, 'et al.' is used.
%
\institute{
	Center of Excellence in Artificial Intelligence, AGH University, al.~Mickiewicza 30, 30-059 Cracow, Poland
	\and
	Nicolaus Copernicus Astronomical Center, Polish Academy of Sciences, ul. Bartycka 18, 00-716 Warsaw, Poland
	\and
	Systems Research Institute, Polish Academy of Sciences, ul.~Newelska 6, 01-447 Warsaw, Poland
}
%
\maketitle              % typeset the header of the contribution
%
\begin{abstract}
	The abstract should briefly summarize the contents of the paper in
	150--250 words.

	\keywords{First keyword  \and Second keyword \and Another keyword.}
\end{abstract}
%
%
%%%%%%%%%%%%%%%%%%%%%%%%%%%%%%%%%%%%%%%%%%%%%%%%%%%%%%%%%%%%%%%%%%%%%%%%%%
\section{Introduction}



%%%%%%%%%%%%%%%%%%%%%%%%%%%%%%%%%%%%%%%%%%%%%%%%%%%%%%%%%%%%%%%%%%%%%%%%%%
\section{Approach}

\begin{equation}
	\langle \sigma_Z^{(j)} \rangle = \bra{\phi(x)}U^\dagger  \sigma_Z^{(j)}  U \ket{\phi(x)}, \qquad j = 1,2,3
\end{equation}

\begin{equation}
	\sigma_Z^{(j)} = \mathbb{I}^{\otimes (j-1)} \otimes \sigma_Z \otimes \mathbb{I}^{\otimes (n-j)}
\end{equation}

\begin{equation}
	\ket{\phi(x)} = \left( R_X(x_1) \otimes R_X(x_2) \otimes ... \otimes R_X(x_n) \right) \ket{0}
\end{equation}

\begin{equation}
	U(\boldsymbol{\theta}) = \prod_{\ell=1}^{L} U_{\mathrm{BEL}}(\boldsymbol{\theta}^{(\ell)})
\end{equation}

\begin{equation}
	U_{\mathrm{BEL}}(\boldsymbol{\theta}) =
	\mathrm{CNOT}_{n,1}
	\left(\prod_{i=1}^{n-1} \mathrm{CNOT}_{i,i+1}\right)
	\left(\bigotimes_{i=1}^{n} R_Y(\theta_i)\right)
\end{equation}

\begin{align*}
	R_Y(\theta) & = e^{-i \theta \sigma_Y / 2}                     \\
	            & = \begin{pmatrix}
		                cos(\frac{\theta}{2}) & -sin(\frac{\theta}{2}) \\
		                sin(\frac{\theta}{2}) & cos(\frac{\theta}{2})
	                \end{pmatrix}
\end{align*}

\begin{equation}
	\mathrm{CNOT}_{i,j}
	=
	\ket{0}\!\bra{0}_i \otimes I_{j}
	+
	\ket{1}\!\bra{1}_i \otimes X_{j}
	, \qquad i \neq j
\end{equation}



%%%%%%%%%%%%%%%%%%%%%%%%%%%%%%%%%%%%%%%%%%%%%%%%%%%%%%%%%%%%%%%%%%%%%%%%%%
\section{Results}

% ===== Hyperparameters: QMLP\_initial\_run =====
\begin{table}[t]
	\centering
	\caption{Hyperparameters for experiment QMLP\_initial\_run}
	\begin{tabular}{ll}
		\hline
		Parameter               & Mean Value                                                                               \\
		\hline
		accelerator             & \texttt{auto}                                                                            \\
		batch\_size             & \texttt{32}                                                                              \\
		callbacks               & \texttt{[<aqml4msc.logging.mlflow\_utils.EpochMetricsTracker object at 0x7525b29f4560>]} \\
		devices                 & \texttt{auto}                                                                            \\
		digits                  & \texttt{[5, 6, 7]}                                                                       \\
		enable\_checkpointing   & \texttt{False}                                                                           \\
		enable\_progress\_bar   & \texttt{True}                                                                            \\
		hidden\_dim             & \texttt{[128]}                                                                           \\
		input\_dim              & \texttt{14}                                                                              \\
		logger                  & \texttt{False}                                                                           \\
		loss\_fn                & \texttt{CrossEntropyLoss()}                                                              \\
		lr                      & \texttt{0.001}                                                                           \\
		max\_epochs             & \texttt{100}                                                                             \\
		model\_name             & \texttt{QMLP\_1}                                                                         \\
		n\_folds                & \texttt{5}                                                                               \\
		n\_layers               & \texttt{3}                                                                               \\
		n\_qubits               & \texttt{6}                                                                               \\
		num\_classes            & \texttt{3}                                                                               \\
		num\_sanity\_val\_steps & \texttt{0}                                                                               \\
		num\_workers            & \texttt{4}                                                                               \\
		parent\_run\_name       & \texttt{QMLP\_initial\_run}                                                              \\
		seed                    & \texttt{42}                                                                              \\
		\hline
	\end{tabular}
\end{table}

% ===== Final Metrics: QMLP\_initial\_run =====
\begin{table}[t]
	\centering
	\caption{Final metrics for QMLP\_initial\_run}
	\begin{tabular}{lc}
		\hline
		Metric              & Value             \\
		\hline
		accuracy\_avg       & $0.884 \pm 0.023$ \\
		accuracy            & $0.886 \pm 0.022$ \\
		f1\_macro           & $0.885 \pm 0.023$ \\
		f1\_micro           & $0.886 \pm 0.022$ \\
		f1\_weighted        & $0.886 \pm 0.023$ \\
		precision\_macro    & $0.886 \pm 0.024$ \\
		precision\_micro    & $0.886 \pm 0.022$ \\
		precision\_weighted & $0.887 \pm 0.023$ \\
		recall\_macro       & $0.884 \pm 0.023$ \\
		recall\_micro       & $0.886 \pm 0.022$ \\
		recall\_weighted    & $0.886 \pm 0.022$ \\
		\hline
	\end{tabular}
\end{table}

% ===== CV Metrics per Fold: QMLP\_initial\_run =====
\begin{table}[t]
	\centering
	\caption{Cross-validation metrics per fold for QMLP\_initial\_run}
	\begin{tabular}{lccccc}
		\hline
		Metric              & Fold 5 & Fold 4 & Fold 3 & Fold 2 & Fold 1 \\
		\hline
		accuracy            & 0.880  & 0.855  & 0.880  & 0.905  & 0.911  \\
		accuracy\_avg       & 0.880  & 0.852  & 0.878  & 0.903  & 0.909  \\
		f1\_macro           & 0.879  & 0.852  & 0.879  & 0.903  & 0.910  \\
		f1\_micro           & 0.880  & 0.855  & 0.880  & 0.905  & 0.911  \\
		f1\_weighted        & 0.880  & 0.855  & 0.880  & 0.905  & 0.911  \\
		precision\_macro    & 0.879  & 0.852  & 0.882  & 0.903  & 0.913  \\
		precision\_micro    & 0.880  & 0.855  & 0.880  & 0.905  & 0.911  \\
		precision\_weighted & 0.881  & 0.854  & 0.881  & 0.905  & 0.912  \\
		recall\_macro       & 0.880  & 0.852  & 0.878  & 0.903  & 0.909  \\
		recall\_micro       & 0.880  & 0.855  & 0.880  & 0.905  & 0.911  \\
		recall\_weighted    & 0.880  & 0.855  & 0.880  & 0.905  & 0.911  \\
		train\_acc          & 0.890  & 0.863  & 0.892  & 0.913  & 0.923  \\
		train\_f1           & 0.890  & 0.862  & 0.892  & 0.913  & 0.923  \\
		train\_loss         & 0.620  & 0.620  & 0.614  & 0.464  & 0.438  \\
		val\_acc            & 0.880  & 0.852  & 0.878  & 0.903  & 0.908  \\
		val\_f1             & 0.880  & 0.852  & 0.879  & 0.903  & 0.909  \\
		val\_loss           & 0.630  & 0.631  & 0.632  & 0.476  & 0.447  \\
		\hline
	\end{tabular}
\end{table}




% ===== Final Metrics: HPO\_classical\_MLP\_2 =====
\begin{table}[t]
	\centering
	\caption{1 random metrics for HPO\_classical\_MLP\_2}
	\begin{tabular}{lc}
		\hline
		Metric              & Value             \\
		\hline
		accuracy\_avg       & $0.925 \pm 0.016$ \\
		accuracy            & $0.925 \pm 0.017$ \\
		f1\_macro           & $0.924 \pm 0.017$ \\
		f1\_micro           & $0.925 \pm 0.017$ \\
		f1\_weighted        & $0.925 \pm 0.017$ \\
		precision\_macro    & $0.927 \pm 0.015$ \\
		precision\_micro    & $0.925 \pm 0.017$ \\
		precision\_weighted & $0.930 \pm 0.013$ \\
		recall\_macro       & $0.925 \pm 0.016$ \\
		recall\_micro       & $0.925 \pm 0.017$ \\
		recall\_weighted    & $0.925 \pm 0.017$ \\
		\hline
	\end{tabular}
\end{table}

%%%%%%%%%%%%%%%%%%%%%%%%%%%%%%%%%%%%%%%%%%%%%%%%%%%%%%%%%%%%%%%%%%%%%%%%%%
\section{Conclusion}







%%%%%%%%%%%%%%%%%%%%%%%%%%%%%%%%%%%%%%%%%%%%%%%%%%%%%%%%%%%%%%%%%%%%%%%%%%
\begin{credits}
	\subsubsection{\ackname} A bold run-in heading in small font size at the end of the paper is
	used for general acknowledgments, for example: This study was funded
	by X (grant number Y).

	\subsubsection{\discintname}
	The authors have no competing interests to declare that are
	relevant to the content of this article.
\end{credits}
%
% ---- Bibliography ----
%
% BibTeX users should specify bibliography style 'splncs04'.
% References will then be sorted and formatted in the correct style.
%
% \bibliographystyle{splncs04}
% \bibliography{mybibliography}
%
\begin{thebibliography}{8}
	\bibitem{ref_article1}
	Author, F.: Article title. Journal \textbf{2}(5), 99--110 (2016)

	\bibitem{ref_lncs1}
	Author, F., Author, S.: Title of a proceedings paper. In: Editor,
	F., Editor, S. (eds.) CONFERENCE 2016, LNCS, vol. 9999, pp. 1--13.
	Springer, Heidelberg (2016). \doi{10.10007/1234567890}


	\bibitem{ref_url1}
	LNCS Homepage, \url{http://www.springer.com/lncs}, last accessed 2023/10/25
\end{thebibliography}






\end{document}
